\section{Introduction}

The simultaneous decline of Moore's law and Dennard scaling, coupled with the rapidly increasing computational demands of modern AI/ML systems, has pushed GPU manufacturers toward increasingly specialized mechanisms for sustaining throughput improvements across generations. Tensor Cores and Matrix Cores have become central to this progression by providing substantially higher matrix multiply-accumulate (MMA) throughput than conventional SIMT floating-point pipelines. NVIDIA Tensor Core FP16 throughput increased by nearly $8\times$ from Volta to Hopper, with further architectural changes introduced in Blackwell~\cite{nvidia2017volta, nvidia2020ampere, nvidia2022hopper, nvidia2024blackwell}. AMD Matrix Cores have similarly evolved across four generations of CDNA architectures~\cite{amd2020cdna1, amd2021cdna2, amd2023cdna3, amd2025cdna4}.

These gains have come not only from adding arithmetic units, but also from larger MMA tiles, wider mixed-precision datapaths, structured sparsity, shared-memory operand delivery, and cooperative multi-warp execution. Tensor Core performance therefore depends on interactions across the SIMT hierarchy: how work is distributed among threads within a warp, how warps cooperate on a larger tile, how Tensor Core blocks scale with issue width, and whether the register file and shared-memory subsystem can sustain the resulting compute throughput.

However, the architectural principles governing this scaling remain poorly understood. Numerous microbenchmarking studies have characterized NVIDIA Tensor Cores~\cite{jia2018dissectingnvidiavoltagpu, 8695642, markidis2018tensorcore, martineau2019v100, jia2019dissectingnvidiaturingt4, yan2020demystifying, 9926299, 9931992, 10579250, luo2025dissectingnvidiahopperarchitecture, 10.1145/3725843.3756041, 11575375} and AMD Matrix Cores~\cite{10590025, kurzynski2025mfma, jarmusch2026microbenchmark}, but commercial GPUs expose only fixed and proprietary implementations. Researchers can measure a particular warp width, MMA tile, or pipeline, but cannot vary the underlying Tensor Core organization, modify operand staging, or synthesize competing designs. Existing studies therefore reveal \emph{what} deployed hardware does, but provide limited insight into \emph{why} a configuration was selected or how it would scale under different constraints.

Studying these questions requires an open and synthesizable GPU platform that exposes the arithmetic array and the surrounding SIMT mechanisms that determine its utilization, including register mapping, operand collection, micro-op scheduling, warp scheduling, shared-memory banking, accumulator capacity, and metadata movement. Simulation-only models may abstract these mechanisms and standalone matrix engines omit their coupling to the SIMT pipeline (Table~\ref{tab:related_work_comparison}). FourShadow studies Tensor Cores as tightly integrated GPU functional units using a configurable, cycle-accurate, and synthesizable implementation built on the open-source RISC-V-based Vortex GPGPU.

% float placed early: renders at the top of page 2
\begin{table*}[t]
\centering
\caption{Comparison of GPU tensor-core architectures, simulators and related research.}
\label{tab:related_work_comparison}
\renewcommand{\arraystretch}{0.9}
\resizebox{\textwidth}{!}{%
\begin{tabular}{p{2.4cm} p{2.4cm} p{4.1cm} p{2.6cm} p{2.2cm} p{1.5cm} c c c}
\toprule
\textbf{Paper} &
\textbf{Simulation} &
\textbf{Data Types} &
\textbf{Thread Config.} &
\textbf{Sparsity} &
\textbf{TC Style} &
\makecell{\textbf{Warp-group}\\\textbf{MMA}} &
\makecell{\textbf{Low Latency}\\\textbf{Dot Product}} &
\makecell{\textbf{Open}\\\textbf{Source}} \\
\midrule
\rowcolor{blue!8}
\textbf{FourShadow (Ours)}
  & Cycle-sim, RTL
  & TF32, FP16/8, BF16/8, INT8/4
  & 2, 4, 8, 16, 32
  & 2:4 structured
  & SIMT
  & \cyes
  & Fused (RTL)
  & \cyes \\
\midrule
\textbf{Virgo}~\cite{kim2025virgo}
  & Cycle-sim, RTL
  & INT8/32
  & Fixed 32
  & \cno
  & Systolic
  & \cyes
  & Discrete
  & \cyes \\
\midrule
\textbf{CoopWarp}~\cite{nada2025cooperative}
  & Cycle-sim
  & FP8, FP16, BF16, FP32
  & 4, 8, 16, 32
  & \cno
  & SIMT
  & \cno
  & Discrete
  & \cyes \\
\midrule
\textbf{RM-STC}~\cite{huang2023rmstc}
  & Cycle-sim
  & FP16, INT8
  & Fixed 32
  & Any degree
  & SIMT
  & \cno
  & Discrete
  & \cno \\
\midrule
\textbf{VEGETA}~\cite{jeong2023vegeta}
  & Cycle-sim, RTL
  & BF16, FP32
  & No SIMT
  & N:M flexible
  & Systolic
  & \cno
  & Discrete
  & \cno \\
\midrule
\textbf{SIMD$^2$}~\cite{zhang2022simd2}
  & HW emulation
  & FP16, INT8
  & Fixed 32
  & \cno
  & SIMT
  & \cno
  & Discrete
  & \cno \\
\midrule
\textbf{Accel-Sim}~\cite{accel-sim}
  & Cycle-sim
  & FP16
  & Fixed 32
  & \cno
  & SIMT
  & \cyes
  & Fused (Sim)
  & \cyes \\
\midrule
\textbf{Sparse TC}~\cite{zhu2019sparsetensorcore}
  & Cycle-sim
  & FP16, INT8
  & Fixed 32
  & Vector N:M
  & SIMT
  & \cno
  & Discrete
  & \cno \\
\midrule
\textbf{GPGPU-Sim}~\cite{raihan2019modeling}
  & Cycle-sim
  & FP16/32
  & Fixed 32
  & \cno
  & Modeled
  & \cno
  & Fused (Sim)
  & \cyes \\
\bottomrule
\addlinespace[2pt]
\multicolumn{9}{l}{Warp-group MMA: several warps share one instruction with shared-memory operands.}
\end{tabular}%
}
\end{table*}

Modern SIMT Tensor Cores are commonly organized as arrays of fused dot-product (FEDP) units rather than as the systolic arrays of standalone dataflow accelerators such as Google TPUs~\cite{jouppi2017indatacenterperformanceanalysistensor, jouppi2020domain, norrie2021design, jouppi2021ten, jouppi2023tpuv4} and AWS Trainium~\cite{aws_trainium_architecture}. A systolic array amortizes operand movement for large, regular matrices, but its rigid dataflow can reduce utilization on small, irregular, or sparse tiles and requires approximately $(2N-1)$ cycles to fill and drain an $N\times N$ array, whereas an FEDP array sits beside the other functional units of a SIMT pipeline and executes larger tiles as micro-operations, which provides low-latency, flexible execution but makes throughput sensitive to the surrounding SIMT organization. Warpgroup MMA expands this design space by allowing multiple warps to cooperatively consume shared-memory operands across several Tensor Core blocks, but more blocks do not guarantee proportional speedup: performance may instead be limited by shared-memory bandwidth, operand-buffer refills, accumulator dependencies, descriptor delivery, micro-op scheduling, or insufficient parallel work. Structured sparsity adds an interaction between the logical matrix representation and the physical datapath: commercial GPUs support 2:4 structured sparsity, in which at least two of every four contiguous elements of matrix A are zero~\cite{nvidia2020ampere, mishra2021acceleratingsparsedeepneural} and the nonzero values are stored in compressed form alongside metadata identifying their positions. Although this nominally halves matrix-A storage and arithmetic work, the realized benefit depends on whether the compressed mapping preserves FEDP utilization (Section~\ref{sec:design_sparse}).

To investigate these interactions, we introduce \textit{\textbf{FourShadow}}, an open-source framework for exploring how SIMT Tensor Core microarchitecture scales across the threads/warp, warps/warpgroup, and issue-width--cores/cluster hierarchy. FourShadow is the white-box counterpart of a commercial tensor core: like a shadow, it outlines a design that cannot be opened, in four components (base tensor core, warp-group tensor core, sparse tensor core, programming model), and we learn from where the two agree and where they contrast. Figure~\ref{fig:figure_integration} illustrates its modifications to Vortex. FourShadow exposes thread count, warpgroup size, issue width, accumulator-register count, operand-staging mode, FEDP dot-product length, descriptor-storage location, arithmetic format, and sparse execution as configurable parameters. We evaluate this design space through microarchitectural characterization, end-to-end applications, FPGA implementation, and ASIC synthesis using ASAP7.

The key contributions of our work can be summarized as follows:

\begin{itemize}
    \item We develop \textit{FourShadow}, an open-source and synthesizable framework for modifying and evaluating SIMT Tensor Core microarchitecture on the RISC-V-based Vortex GPGPU.

    \item We evaluate Tensor Core scaling across threads/warp, warps/warpgroup, and issue width--cores/cluster, demonstrating that arithmetic resources alone provide diminishing returns as operand delivery and non-TCU execution become dominant and that the best configuration, including the preferred warp width, depends on workload behavior rather than arithmetic throughput alone: warp width, warpgroup size, accumulator capacity, and FEDP count must be co-designed with operand delivery, register pressure, shared-memory bandwidth, sparsity mapping, and available workload parallelism.

    \item We design a configurable warpgroup MMA architecture supporting register--shared-memory (RS) and shared-memory--shared-memory (SS) operand staging, shared B-tile broadcast, private per-warp A buffering, and configurable accumulator capacity, and analyze how operand delivery, micro-op scheduling, accumulator reuse, and shared-memory organization affect MMA throughput and operational intensity. We explore dot-product scaling through a 2K-wide FEDP design and evaluate its throughput--pipeline-latency tradeoff, finding that wider dot-product units remain beneficial despite deeper pipelines. We also study moving WGMMA descriptors into TCU-local SRAM, showing that its memory, synchronization, and AGU overhead outweighs the reduction in register pressure.

    \item We develop a complete 2:4 structured-sparsity datapath and show that sparse acceleration depends strongly on symmetric versus asymmetric Tensor Core configurations and RS versus SS operand staging, with sparse WGMMA providing $2.21$--$2.35\times$ speedup when compressed mappings preserve FEDP utilization.

    \item We validate FourShadow through cycle-accurate simulation, FPGA implementation, and ASIC synthesis. On an AMD Alveo U55C FPGA at 250\,MHz, WGMMA increases end-to-end speedup over the SIMT baseline from $5.91\times$ to $9.15\times$ for Llama~2 prefill and from $1.42\times$ to $2.00\times$ for decode, while the combined extensions reach $16.88\times$ on CNN and $26.82\times$ on NeRF workloads.
\end{itemize}